\documentclass[12pt, letterpaper]{article}

%-------Packages---------
\usepackage{newpxtext,newpxmath} % Palatino-like text & math (modern)
\usepackage{bboldx}
\usepackage{mathtools}
\usepackage{tikz}
\usetikzlibrary{calc,intersections}
\usepackage{tikz-cd}
\usepackage[all,arc]{xy}
\usepackage{graphicx}

\usepackage{enumerate}
\usepackage[dvipsnames]{xcolor}
\usepackage{float}
\usepackage[left=1in,top=1in,right=1in,bottom=1in]{geometry}
\usepackage{fancyhdr}
\usepackage{multicol}
\usepackage{setspace}

\usepackage{dsfont}
\usepackage{mathrsfs}

\usepackage{tcolorbox}
\tcbuselibrary{theorems,skins,breakable}

\usepackage{xparse}
\usepackage{import}
\usepackage[utf8]{inputenc}

\usepackage{hyperref}
\usepackage{cleveref}

\usepackage{xcolor, ifthen, tcolorbox}
\tcbuselibrary{theorems,skins,breakable}
% Theorem System
% The following environments are provided:
%   New Problem:        \begin{problem} ...
%   New Question Header:   \qh (then followed by subproblems)
%   Subproblem:     \begin{subproblem} ...
%   Problem Proof:  \begin{problemproof} ...
%   Proof:          \\begin{pf}{color} ...
%   Remark:         \rmk (sentence), \rmkb (block)

% Define custom colors
\definecolor{thmscol}{HTML}{F4565B} %For Problems
\definecolor{lemscol}{HTML}{FE844C} %For Lemmes
\definecolor{prpscol}{HTML}{00A7EF} %For proposition
\definecolor{corscol}{HTML}{23DAFF} %For corrolaries
\definecolor{rmkscol}{HTML}{6DEEC5} %For remarks
\definecolor{defscol}{HTML}{18C991} %For definitions
\definecolor{exmscol}{HTML}{7DB800} %For examples
\definecolor{clmscol}{HTML}{165c58} %For claims

%Change between dark(black)/light(white) mode
\newcommand{\colormode}{white}
\ifthenelse{\equal{\colormode}{white}}
      {\newcommand{\TextColor}{black}}
      {\newcommand{\TextColor}{white}}
\newcommand{\pbcolormain}{thmscol!80!\colormode}
\newcommand{\pbcolorsecond}{thmscol!38!\colormode}


% ============================
% Problem Counter
% ============================
\newcounter{subproblemcounter}
\newcounter{problemcounter}

\newcommand{\q}{
    \setcounter{subproblemcounter}{0}%
    \stepcounter{problemcounter} % Increment problem counter
    \color{\TextColor}Problem \arabic{problemcounter}: % Display the problem counter
}

\newcommand{\stepsub}{
    \stepcounter{subproblemcounter}%
    \color{\TextColor}(\alph{subproblemcounter})
}
% ============================
% Problem
% ============================
\newtcolorbox{pb}[1]{
  enhanced,
  frame hidden,
  titlerule=0mm,
  toptitle=1mm,
  bottomtitle=1mm,
  fonttitle=\bfseries\large,
  coltitle=black,
  colbacktitle=\pbcolormain,
  colback=\pbcolorsecond,
  title={\q #1},
  coltext=\TextColor
}

\newtcolorbox{spb}[1]{
  enhanced,
  frame hidden,
  titlerule=0mm,
  toptitle=1mm,
  bottomtitle=1mm,
  fonttitle=\bfseries\large,
  coltitle=black,
  colbacktitle=\pbcolormain,
  colback=\pbcolorsecond,
  title={\stepsub #1},
  coltext=\TextColor,
  attach boxed title to top left={xshift=3mm,yshift=-2mm},
  boxed title style={
    colback=\pbcolormain,
    colframe=\pbcolormain,
  }
}

\newtcolorbox{pbh}[1]{
    enhanced,
    frame hidden,
    titlerule=0mm,
    toptitle=1mm,
    bottomtitle=1mm,
    fonttitle=\bfseries\large,
    coltitle=black,
    colbacktitle=\pbcolormain,
    colback=\pbcolorsecond,
    title={\q #1},
    boxrule=0pt,
    interior hidden,
    attach boxed title to top left={xshift=3mm,yshift=-2mm},
    boxed title style={
        colback=\pbcolormain,
        colframe=\pbcolormain,
    }
}

\newenvironment{problem}[1]{%
  \newpage
  \begin{pb}{#1}
    }{
  \end{pb}
}

\newenvironment{subproblem}[1]{%
  \begin{spb}{#1}
    }{
  \end{spb}
}

\newenvironment{problemproof}{%
    \begin{pf}{\pbcolormain}
        }{
    \end{pf}
}

\newcommand{\qh}[1][]{
    \newpage
    \begin{pbh}{#1}\end{pbh} % Display the problem counter
}

% ============================
% Claim
% ============================

\newtcbtheorem[number within=problemcounter]{claim}{Claim}
{
    enhanced,
    frame hidden,
    titlerule=0mm,
    toptitle=1mm,
    bottomtitle=1mm,
    fonttitle=\bfseries\large,
    coltitle=black,
    colbacktitle=clmscol!50!\colormode,
    colback=clmscol!20!\colormode,
}{clm}

\NewDocumentCommand{\clm}{m+m}{
    \begin{claim*}{#1}{}
        #2
    \end{claim*}
}

\newenvironment{clmpf}{
	{\noindent{\it \textbf{Proof.}}
	\setlength{\parindent}{0pt}
	}
	\tcolorbox[blanker,breakable,left=5mm,parbox=false,
    before upper={\parindent15pt},
    after skip=10pt,
	borderline west={1mm}{0pt}{clmscol!50!\colormode}]
}{
    \textcolor{clmscol!50!\colormode}{\hbox{}\nobreak\hfill$\blacksquare$}
    \endtcolorbox
}

\NewDocumentCommand{\clmp}{m+m+m}{
    \begin{claim*}{#1}{}
        #2
    \end{claim*}

    \begin{clmpf}
        #3
    \end{clmpf}
}


% ============================
% Proof
% ============================

\newenvironment{pf}[1]{%
  \def\pfcolor{#1}% Store the color in a macro
  \noindent{\it \textbf{Proof.}}%
  \tcolorbox[blanker,breakable,left=5mm,parbox=false,%
    before upper={\parindent15pt},%
    after skip=10pt,%
    borderline west={1mm}{0pt}{\pfcolor},%
    coltext=\TextColor
  ]%
  \setlength{\parindent}{0pt}%
}{%
  \textcolor{\pfcolor}{\hbox{}\nobreak\hfill$\blacksquare$}%
  \endtcolorbox%
}

\newtcolorbox{solution}{
  blanker,
  breakable,
  left=5mm,
  parbox=false,%
  before upper={\parindent15pt},%
  after skip=10pt,%
  borderline west={1mm}{0pt}{\pbcolormain},%
  coltext=\TextColor
}


% ============================
% Remark
% ============================


\NewDocumentCommand{\rmk}{+m}{
    {\it \color{rmkscol!100!white}#1}
}

\newenvironment{remark}{
    \par
    \vspace{5pt}
    \begin{minipage}{\textwidth}
        {\par\noindent{\textbf{Remark.}}}
        \tcolorbox[blanker,breakable,left=5mm,
        before skip=10pt,after skip=10pt,
        borderline west={1mm}{0pt}{rmkscol!80!\colormode},
        coltext=\TextColor
        ]
}{
        \endtcolorbox
    \end{minipage}
    \vspace{5pt}
}

\NewDocumentCommand{\rmkb}{+m}{
    \begin{remark}
        #1
    \end{remark}
}


% Define a new command for problems
\renewcommand{\headrule}{\color{\TextColor}\hrulefill}
\newcommand{\C}{\mathbb{C}}
\newcommand{\R}{\mathbb{R}}
\newcommand{\Q}{\mathbb{Q}}
\newcommand{\Z}{\mathbb{Z}}
\newcommand{\N}{\mathbb{N}}
\newcommand{\p}{\mathfrak{p}}
\newcommand{\F}{\mathbb{F}}
\newcommand{\contr}{\Rightarrow \Leftarrow}
\newcommand{\s}{\(\,\)}
\renewcommand{\t}[1]{\text{#1}}
\newcommand{\Spec}{\mathrm{Spec}}
\newcommand{\Ass}{\mathrm{Ass}}
\newcommand{\Supp}{\mathrm{Supp}}
\newcommand{\Ann}{\mathrm{Ann}}
\newcommand{\mf}[1]{\mathfrak{#1}}

% Meta data
\setstretch{1.2}

\begin{document}
\color{\TextColor}
\pagecolor{\colormode}
\setlength{\parindent}{0pt}
\pagestyle{fancy}
\fancyhf{}
\fancyhead[LOE]{\color{\TextColor}\slshape{\textbf{Vincent Ferrara}}}
\fancyhead[ROE]{\color{\TextColor}HW9 --- \thepage}
\fancyhead[C]{\color{\TextColor}Math 614}


\begin{problem}{}
    Let $R$ be a noetherian ring such that $\Spec(R)$ is finite. Prove that $\dim R \leq 1$, and give an
example for which $\dim R = 1$.
\end{problem}
\begin{problemproof}
    Assume $R$ is noetherian and that $\Spec(R)$ is finite.

    Let $\p \in \Spec(R)$. Let $\{\mf{q}_1,\dots,\mf{q}_n\}$ be the primes in $\Spec(R)$ that do not contain $\p$.
    This is a finite list since $\Spec(R)$ is finite.

    By the prime avoidence lemma this implies that $\p \nsubset \bigcup_{i=1}^n \mf{q}_i$.

    This implies there exists an $x \in \p \setminus \bigcup_{i=1}^n \mf{q}_i$.

    If $\tau \nsupset \p$ and $\tau \in \Spec(R)$ this implies that $\tau = \mf{q}_i$ for some $i$. Thus, $x \notin \tau$.
    
    If $\tau \supset \p$ and $\tau \in \Spec(R)$ this implies that $x \in \tau$.

    Thus, $x \in \tau \iff \p \subset \tau$.

    This implies $V(x) = \{\mf{q} \in \Spec(R) \mid x \in \mf{q}\} = \{\mf{q} \in \Spec(R) \mid \p \subset \mf{q}\} = V(\p)$.

    This implies that $\sqrt{(x)} = \sqrt{\p} = \p$.

    Thus, $\p$ is minimal over $(x)$.

    Therefore, by the principle ideal theorem this implies that $\dim \p \leq 1$.

    Thus, since $\p$ is arbitrary we get that $\dim R \leq 1$.

    Example $\Z_{(p)}$ s.t. $p$ is prime. The prime ideals in here are $(0),(p)R_{(p)}$. Thus, $\dim \Z_{(p)} = 1$.
\end{problemproof}

\begin{problem}{}
    Let $R$ be a noetherian ring, and let $R[x]$ be the polynomial ring over $R$ in one variable. Let
$R[x, x^{-1}]$ be the localization of $R[x]$ at the element $x$. Prove that $\dim R[x, x^{-1}] = \dim R + 1$.

\end{problem}
\begin{problemproof}
    Assume $R$ is noetherian.

    Shown in class $\dim R[x] = \dim R + 1$. Also, shown in class since $R[x,x^{-1}]$ is a localization of $R[x]$ this implies $\dim R[x,x^{-1}] \leq \dim R[x] = \dim R + 1$.

    First if $\dim R = \infty$ then we can extend a chain to $R[x]$ by just changing $\p$ with $\p[x]$ which was proven prime in class.
    Then extending $R[x,x^{-1}]$ since $x \notin \p[x]$ for any prime $\p$ this implies that this is also a chain in $R[x,x^{-1}]$ since primes in $R[x]$ s.t. $x$ is not in the prime are in bijection with primes in $R[x,x^{-1}]$.

    Thus, $\dim R[x,x^{-1}] = \infty = \dim R + 1$

    If $\dim R$ is finite.
    
    Let $\p_0 \subsetneq \p_1 \subsetneq \dots \subsetneq \p_n$ be a max chain of prime ideals in $R$ where $n = \dim R$.

    First we can extend this to $\p_0[x] \subsetneq \p_1[x] \subsetneq \dots \subsetneq \p_n[x]$.
    
    Also consider $\p_n[x] + (x+1)$ this is also prime since $R[x] \diagup (\p_n[x] + (x+1)) \cong R \diagup \p_n$.

    Thus, $\p_0[x] \subsetneq \p_1[x] \subsetneq \dots \subsetneq \p_n[x] \subsetneq \p_n[x] + (x+1)$ is a chain. Also, since $x$ is not in any of these prime ideals.
    
    This implies that this is also a chain in $R[x,x^{-1}]$. Thus, $\dim R[x,x^{-1}] \geq \dim R + 1$.

    $\dim R[x,x^{-1}] = \dim R + 1$.
\end{problemproof}

\begin{problem}{}
    Let $R$ be a noetherian ring.
\end{problem}
\begin{subproblem}{}
    Show that if $\p \subset R$ is a prime ideal, then for any $n \geq 1$ its $n$th symbolic power $\p^{(n)}$
is the smallest $\p$-primary ideal containing $\p^n$. That is, if $\mf{q}$ is another $\p$-primary ideal
containing $\p^n$, then $\p^{(n)} \subset \mf{q}$.
\end{subproblem}
\begin{problemproof}
    Let $\p$ be a prime ideal of $R$.

    Assume $\mf{q}$ is $\p$-primary ideal s.t. $\p^n \subset \mf{q}$.

    Let $x \in \p^{(n)} \implies \exists s \notin \p$ s.t. $sx \in \p^n \subset \mf{q}$.

    Thus, $xs \in \mf{q}$. Since $\mf{q}$ is $\p$-primary this implies either $x \in \mf{q}$ or $s^n \in \p$, but if $s^n \in \p \implies s \in \p$ which is a contradiction.

    Therefore, $x \in \mf{q}$. Thus, $\p^{(n)} \subset \mf{q}$.
\end{problemproof}

\begin{subproblem}{}
    Show that if $\mf{m} \subset R$ is a maximal ideal, then $m^n = m^{(n)}$ for all $n \geq 1$
\end{subproblem}
\begin{problemproof}
    Let $\mf{m}$ is a maximal ideal.

    Consider $\mf{m}^n$. First $\mf{m}^n \subset \mf{m}^{(n)}$
    
    Since $\sqrt{\mf{m}^n}=\mf{m}$ is maximal this implies that $\mf{m}^n$ is $\mf{m}$-primary.
    
    Thus, by $(a)$ this implies that $\mf{m}^{(n)} \subset \mf{m}^n$.

    Therefore, $\mf{m}^{(n)} = \mf{m}^n$
\end{problemproof}

\begin{problem}{}
    Let $R$ be a noetherian ring and $f \in R$.
\end{problem}
\begin{subproblem}{}
    Prove that if $f$ is non-zerodivisor in $R$, then $\dim R \diagup (f) \leq \dim R -1$.
\end{subproblem}
\begin{problemproof}
    Assume $f$ is a non-zerodivisor.

    Assume for contradiction that $f \in \p$ s.t. $\p$ is a minimal prime.

    Thus, consider $R_\p$. Since there is only one maximal ideal in the localization this implies it is artinian.

    Thus, $\p R_\p^n =0$ for some $n \in \N$. This implies that $f^n/1 =0$. Thus, there exists an $s \notin \p$ s.t. $sf^n=0$. Since $f$ is a non-zero divisor this implies $s=0$. This is a contradiction though.

    Thus, $f$ is not in any minimal prime.

    Therefore, the length of chains of primes that contain $f$ is at most length $\dim R - 1$ since $f$ is not in any minimal primes.

    Therefore, since there is a bijection between primes in $R$ that contain $f$ and primes in $R \diagup (f)$. This implies that $\dim R \diagup (f) \leq \dim R - 1$.
\end{problemproof}

\begin{subproblem}{}
    If $(R, \mf{m})$ is a local ring and $f \in \mf{m}$, prove that $\dim R-1 \leq \dim R \diagup (f) \leq \dim R$.
In particular, if $f$ is also a non-zerodivisor, then by the previous part we have
$\dim R \diagup (f) = \dim R-1$.
\end{subproblem}
\begin{problemproof}
    Assume $(R, \mf{m})$ is a local ring and $f \in \mf{m}$.

    Let $\p$ be a minimal prime over $f$.

    By the principle ideal theorem this implies that $\mathrm{codim}(\p) \leq 1$.

    This implies a max length chain that starts at $(\p)$ is $\mathrm{codim}(\mf{m})-\mathrm{codim}(\mf{\p}) \geq \dim R - 1$.

    Thus, $\dim R \diagup (f) \geq \dim R - 1$.
\end{problemproof}

\begin{subproblem}{}
    Give an example of a local ring $(R, m)$ and $f \in \mf{m}$ such that $\dim R - 1 = \dim R \diagup (f)$,
and give an example so that $\dim R = \dim R \diagup (f )$.
\end{subproblem}
\begin{problemproof}
    \begin{itemize}
        \item Consider $k[[x]]$ formal power series of a field $k$ in one variable. Proven in 494 $f \in k[[x]]$ is a unit if and only if $a_0 \neq 0$.
        Thus, the set of non units is exactly $(x)$. A proper ideal in $k[[x]]$ cannot contain a unit, thus every proper ideal is contained in $(x)$. Thus, $(x)$ is the unique maximal ideal.

        Let $f \in k[[x]]$ be non zero. Let $a_n$ be the first term that is non zero in $f$. This implies that $f=t^n \cdot u$ s.t. $u$ is a unit since its first term is non-zero which implies it is a unit.

        Let $I \neq 0$ be an ideal. Let $f \in I$ s.t. $f=x^n \cdot u$ where $n$ is minimal amount all elements in $I$. This comes from the well ordering principle since these powers are bounded below by zero and all natural numbers.

        Thus, $x^n = f \cdot u^{-1} \in I$. Thus, $(x^n) \subset I$.
        Let $g \in I$. Write $g=x^m \cdot v$. Since $n$ is minimal we have that $n \leq m$, so $g \in (x^n)$.

        Thus, $I=(x^n)$. Thus, every non-zero ideal in $k[[x]]$ is of the form $(x^n)$ for $n \geq 0$.

        Lets consider prime ideals now
       
        Since $k[[x]]$ is a domain $(0)$ is prime.

        Now let $\p \neq 0$ be prime. Let $0 \neq f \in \p$. As above we have that $f=x^n u$ since $u$ is a unit this implies $x^n \in \p \implies x \in \p$. Thus, $x \in \p$, so $\p = (x)$. Therefore,
        the only prime ideals are $(0)$ and $(x)$. Thus, the max chain is $(0) \subsetneq (x)$. Thus, $k[[x]]$ is dimension one.

        Also, $k[[x]] \diagup (x) \cong k$ which is dimension zero. Thus, $\dim k[[x]] - 1 = \dim k[[x]] \diagup (x)$.
        \item Consider $R=\R$ then since $(0)$ is the only prime ideal thus maximal we get that $\dim \R = 0 = \dim \R \diagup (0)$
    \end{itemize}
\end{problemproof}

\begin{problem}{}
    
\end{problem}
\begin{problemproof}
    Since $A=k[x]_{(x)}$ this implies that $\dim A \leq \dim k[x]=1$, but since we have the chain $(0) \subsetneq (x)$ in $A$ this implies $\dim A =1$.

    Let $R=A[t]$. Since $A$ is noetherian this implies that $\dim R = \dim A[t] = \dim A + 1 = 2$.

    Consider $A[t] \diagup (xt-1) \cong A[1/x]$.

    Since $A=k[x]_{(x)}$. This implies we inverted everything that has a non zero constant term. Then in $A[1/x]$ we then localize again and invert $x$ which then inverts non-zero polynomials with zero in the constant term.

    Thus, we have inverted $k[x] \setminus {0}$. Therefore, $A[1/x] \cong k(x)$. Therefore, $(xt-1)$ is a maximal ideal.

    Thus, $\dim R \diagup \p = 0$.

    Also, since $R$ is noetherian by the principle ideal theroem we get that $\mathrm{codim}(\p) \leq 1$. Since $R$ is a domain this implies that $(0) \subset \p$. Thus, $\p$ has codimension 1.
\end{problemproof}
\end{document}